\chapter{Sistema de Tipus}

\section{Fonaments de l'Algorisme de Hindley-Milner}

El sistema de tipus d'aquest intèrpret s'ha dissenyat sota el marc formal de l'algorisme de \textbf{Hindley-Milner (HM)}, un sistema de deducció lògica adaptat per a llenguatges funcionals amb polimorfisme paramètric. Com descriu l'obra de referència de Peyton Jones \cite{peyton1987implementation}, la potència d'aquest enfocament rau en el fet que el programador gaudeix de la robustesa d'un llenguatge de tipat estàtic fort sense patir la càrrega sintàctica de realitzar anotacions de tipus explícites a cada declaració.

L'algorisme opera sota un principi d'obvencions axiomàtiques i deduccions estructurals. El sistema calcula formalment el \textit{principal type} (el tipus més general possible), basant-se en les teories formulades originalment per Luis Damas i Robin Milner \cite{damas1982principal}. Aquest rigor matemàtic garanteix el teorema de seguretat fonamental formulat pel mateix Milner: \textit{``Well-typed programs cannot go wrong''}. En el context d'aquest projecte, que un programa estigui ``ben tipat'' significa que el motor d'inferència codificat a \texttt{HSInferenciaTipus.hs} és capaç de trobar una assignació de tipus consistent que respecti les regles sintàctiques primitives. Si l'algorisme detecta una incongruència semàntica, el procés d'execució s'interromp immediatament abans de la fase de reducció de grafs, protegint la integritat del sistema de memòria.

La sintaxi dels tipus suportats es modela de forma abstracta en el fitxer \texttt{HSTipus.hs} mitjançant el tipus de dades algebraic \texttt{Tipus}:

\begin{lstlisting}[style=haskellstyle, caption=Definició formal dels tipus de dades de l'intèrpret]
data Tipus = TInt
           | TBool
           | TFun Tipus Tipus
           | TVar String
           deriving (Show, Eq)
\end{lstlisting}

Aquest model reflecteix fidelment els elements atòmics del llenguatge: els tipus base concrets (\texttt{TInt} per a literals numèrics i \texttt{TBool} per a valors booleans), el constructor de funcions d'ordre superior (\texttt{TFun}), i les variables de tipus abstractes (\texttt{TVar}), que actuen com a marcadors de posició universals durant les fases d'unificació.

\section{L'Algorisme d'Unificació de Robinson i el Test d'Ocurrència}

L'engranatge central que permet resoldre les equacions de tipus generades per l'intèrpret és l'algorisme d'unificació de Robinson. Donats dos tipus qualssevol $\tau_1$ i $\tau_2$, el propòsit d'aquest algorisme és trobar un conjunt de mapatges —anomenat Substitució ($\sigma$)— que aconsegueixi que ambdós tipus esdevinguin estructuralment idèntics quan se'ls aplica la dita transformació ($\sigma(\tau_1) = \sigma(\tau_2)$).

A la nostra implementació, aquest comportament es troba encapsulat en la funció \texttt{generaSubst}. Aquesta funció opera comparant la geometria d'ambdues estructures de dades, resolent dependències creuades i evitant cicles infinits mitjançant un rigorós \textit{Occurs Check} (Test d'Ocurrència):

\begin{lstlisting}[style=haskellstyle, caption=Implementació del motor d'unificació de Robinson amb protecció contra tipus infinits]
generaSubst :: Tipus -> Tipus -> Either String Subst
generaSubst TInt TBool = Left "Error d'inferència: No es pot unificar un enter amb un booleà"
generaSubst TBool TInt = Left "Error d'inferència: No es pot unificar un booleà amb un enter"

generaSubst (TFun t1 t2) (TFun t3 t4) = do 
  first  <- generaSubst t1 t3
  second <- generaSubst (aplicaSubst first t2) (aplicaSubst first t4)
  return (second ++ first) -- Posem les més recents primer

generaSubst (TVar ('t':s1)) (TVar s2) = Right [(('t':s1), (TVar s2))]  
generaSubst (TVar s1) (TVar s2) = Right [(s2, (TVar s1))]  

generaSubst t (TVar s) 
  | TVar s == t = return []
  | ocorreEn s t = Left "Error d'inferència: Tipus infinit"
  | otherwise = return [(s,t)]

generaSubst (TVar s) t = generaSubst t (TVar s)
generaSubst _ _ =  Right []

ocorreEn :: String -> Tipus -> Bool
ocorreEn s (TVar a) = s == a 
ocorreEn s (TFun t1 t2) = ocorreEn s t1 || ocorreEn s t2 
ocorreEn _ _ = False
\end{lstlisting}

L'algorisme es divideix en els següents escenaris operacionals clars:

\begin{enumerate}
    \item \textbf{Conflicte de Tipus Primitius:} L'intent directe d'equiparar \texttt{TInt} amb \texttt{TBool} (o viceversa) representa una violació semàntica irreconciliable. La funció retorna una fallida monàdica \texttt{Left}, propagant un missatge clar d'error cap a la terminal de l'usuari.
    \item \textbf{Descomposició d'Estructures (\texttt{TFun}):} Quan es comparen dues funcions, l'algorisme es ramifica recursivament. És de vital importància notar com primer s'unifiquen els dominis (\texttt{t1} i \texttt{t3}) generant la substitució \texttt{first}. A continuació, per mantenir la coherència lògica, \textbf{s'aplica immediatament} aquesta substitució \texttt{first} als codominis (\texttt{t2} i \texttt{t4}) abans d'intentar unificar-los. Això garanteix que qualsevol resolució de variables feta en els arguments es propagui cap als resultats.
    \item \textbf{El Test d'Ocurrència (\textit{Occurs Check}):} Un dels reptes teòrics més complexos de la inferència és la gestió d'expressions divergents com \texttt{\textbackslash x -> x x}. Intentar inferir això implicaria dir que el tipus de \texttt{x} ha de ser igual a \texttt{x -> y}. Això genera un tipus infinit: \texttt{(... -> y) -> y}. Per evitar que l'unificador entri en un bucle infinit penjant l'intèrpret, s'ha dissenyat la funció auxiliar \texttt{ocorreEn}. Aquesta funció fa un escaneig profund de l'arbre: si s'intenta lligar una variable \texttt{s} a un tipus \texttt{t} que \textit{ja conté} la pròpia variable \texttt{s} en el seu interior, el sistema llança l'error ``Tipus infinit'', avortant la operació de manera segura.
\end{enumerate}


\section{El Motor d'Inferència Estructural}

La deducció de tipus s'executa mitjançant un recorregut sintàctic recursiu sobre l'AST utilitzant la funció \texttt{infereix}. En aquesta implementació evolucionada, la funció requereix la gestió simultània de dos contextos: el context inicial global (\texttt{inictx}), on viuen les primitives del \textit{Prelude}, i el context lèxic local (\texttt{ctx}), que muta a mesura que s'aprofundeix en les expressions. Els contextos es modelen com a llistes de parells: \texttt{type Context = [(String, Tipus)]}.

La propagació d'informació es realitza mitjançant la composició i l'aplicació seqüencial de substitucions. Les funcions auxiliars \texttt{aplicaSubst} i \texttt{aplicaSubstCtx} s'encarreguen de recórrer un tipus o un context sencer per substituir-ne les variables lliures per les noves definicions estructurals descobertes en temps de computació:

\begin{lstlisting}[style=haskellstyle, caption=Mecanisme de propagació de substitucions en tipus i contexts]
aplicaSubst :: Subst -> Tipus -> Tipus
aplicaSubst s (TVar v) = case lookup v s of
                            Just t  -> aplicaSubst s t -- Continuem buscant
                            Nothing -> TVar v
aplicaSubst s (TFun t1 t2) = TFun (aplicaSubst s t1) (aplicaSubst s t2)
aplicaSubst _ t = t

aplicaSubstCtx :: Subst -> Context -> Context
aplicaSubstCtx _ [] = []
aplicaSubstCtx subs ((st,t):ctx) = ((st, aplicaSubst subs t): (aplicaSubstCtx subs ctx)) 
\end{lstlisting}

La funció \texttt{aplicaSubst} realitza una cerca recursiva profunda (\textit{deep lookup}): si una variable \texttt{v} es mapeja cap a una altra variable de tipus, l'algorisme continua explorant la llista de substitucions iterativament per assegurar-se que es resol el tipus terminal més madur disponible a la memòria.

\subsection{Polimorfisme Paramètric i Instanciació de Variables Fresques}

Un dels assoliments més destacats del disseny d'aquest motor és la implementació de polimorfisme paramètric (\textit{Let-Polymorphism}) per a funcions predefinides \cite{peyton1987implementation}. Quan es fa una crida a una funció polimòrfica emmagatzemada al \texttt{preludeTypeEnv} (com ara \texttt{id :: a -> a}), no podem utilitzar la variable \texttt{"a"} directament cada vegada. Si ho féssim, intentar avaluar \texttt{if (id True) then (id 42) else 0} faria que l'algorisme unifiqués \texttt{"a"} primer amb \texttt{Bool} i després, de manera contradictòria, amb \texttt{Int}, llançant un error fatal.

Per solucionar això, cada vegada que s'extreu un tipus del context, s'ha de \textbf{instanciar}: cal substituir les seves variables genèriques originals per "variables fresques" noves, úniques per a aquella crida en concret. Això s'ha solucionat mitjançant una enginyosa estratègia d'acumulació de números reflectida en les funcions \texttt{renameTVar} i \texttt{forceRenameTVar}:

\begin{lstlisting}[style=haskellstyle, caption=Mètodes per a l'instanciació de variables i reanomenament dinàmic]
renameTVar :: Tipus -> Int -> Maybe Tipus
renameTVar (TFun t1 t2) n = do 
  r1 <- (renameTVar t1 n)
  r2 <- (renameTVar t2 n)
  return (TFun r1 r2)
renameTVar (TVar ('t':rest)) n = Just (TVar ('t':rest))
renameTVar (TVar s) n = Just (TVar (s++(show n)))
renameTVar t _ = Nothing

forceRenameTVar :: Tipus -> Int -> Maybe Tipus
forceRenameTVar (TFun t1 t2) n = do 
  r1 <- (forceRenameTVar t1 n)
  r2 <- (forceRenameTVar t2 n)
  return (TFun r1 r2)
forceRenameTVar (TVar s) n = Just (TVar (s++(show n)))
forceRenameTVar t _ = Nothing
\end{lstlisting}

El mecanisme fa servir un comptador enter global (\texttt{n}) que viatja a través de la mònada. Quan un tipus genèric és sol·licitat, aquestes funcions recorren el tipus de dades algebraic. Si troben una variable com \texttt{TVar "a"}, hi concatenen el comptador vigent, transformant-la en \texttt{TVar "a1"}. Així, crides successives generaran signatures aïllades com \texttt{a1 -> a1} o \texttt{a2 -> a2}, prevenint qualsevol col·lisió o pol·lució de l'espai de noms de tipus de la memòria de l'intèrpret. Cal destacar que \texttt{renameTVar} fa una excepció lògica per a les variables que comencen per \texttt{'t'}, ja que aquestes són reconegudes com a variables internes de treball de l'algorisme en curs i no s'han de duplicar, en contrast amb \texttt{forceRenameTVar} que força l'agrupació sense distincions per a lligams preexistents.

\subsection{Anàlisi de Casos Claus de la funció infereix}

L'arquitectura funcional de la deducció estructural s'implementa bifurcant el patró sintàctic de l'expressió analitzada. Els blocs monàdics \texttt{do} gestionen de manera nativa les fallides d'unificació i l'enllaç de la mònada interna de Haskell \cite{wadler1992essence}.

\subsubsection{Cerca de Variables (Variables Lliures i Lligades)}

Quan l'algorisme analitza un node \texttt{Var a}, fa una exploració de presència. Primerament busca la variable en l'entorn local (\texttt{ctx}); si hi és, significa que és una variable capturada localment (per exemple el paràmetre d'una lambda externa). Si no es troba, procedeix a cercar-la a l'entorn global (\texttt{inictx}). És precisament aquí on fa ús de les funcions descrites anteriorment per a instanciar el seu polimorfisme de forma segura, modificant adequadament el comptador d'adreces \texttt{n}.

\subsubsection{Abstraccions Lambda i Paràmetres Pattern-Matched}

El tractament de les expressions \texttt{Lam s e} presenta una dualitat fonamental segons la naturalesa del paràmetre d'entrada:

\begin{lstlisting}[style=haskellstyle, caption=Inferència de tipus per a abstraccions lambda i literals en línia]
infereix inictx ctx (Lam s e) n = do
  case all isDigit s of
    True -> (do  
      t_s <- Right (TInt)                                     
      (t1, s1, n1) <- infereix inictx ctx e n                 
      t_retorn <- Right (aplicaSubst s1 (TFun t_s t1))        
      return (t_retorn, s1, n1)
      )
    False -> (do  
      t_s <- Right (TVar ("t"++ show n ))                     
      (t1, s1, n1) <- infereix inictx ((s,t_s):ctx) e (n+1)   
      t_retorn <- Right (aplicaSubst s1 (TFun t_s t1))        
      return (t_retorn, s1, n1)
      )
\end{lstlisting}

Per donar suport a cert nivell de reconeixement de patrons (\textit{pattern matching}) directament als paràmetres (per exemple \texttt{\textbackslash 1 -> True}), s'utilitza una comprovació \texttt{all isDigit s}. Si el paràmetre és un numeral, l'intèrpret assigna directament el tipus fix \texttt{TInt}. Per contra, el cas general (quan s'empren lletres com \texttt{x} o \texttt{y}) es tracta enllaçant el nou identificador sintàctic amb una variable de tipus de treball fresca \texttt{TVar ("t" ++ show n)}. Aquest supòsit s'injecta a l'inici del context abans de procedir a inferir el cos de l'expressió (\texttt{e}). Finalment, reconstrueix la definició pròpia de la funció mitjançant el constructor \texttt{TFun}.

\subsubsection{Aplicació de Funcions}

L'aplicació d'una funció a un argument (\texttt{App e1 e2}) és el punt on s'exigeix la màxima consistència estructural, ja que s'està instanciant efectivament el còmput:

\begin{lstlisting}[style=haskellstyle, caption=Inferència i unificació acumulativa de l'aplicació de funcions]
infereix inictx ctx (App e1 e2) n = do
  (t1, s1, n1) <- infereix inictx ctx e1 n                       
  (t2, s2, n2) <- infereix inictx (aplicaSubstCtx s1 ctx) e2 n1  
  
  t3 <- Right (TVar ("t" ++ show n2))                     
  s3 <- generaSubst (aplicaSubst s2 t1) (TFun t2 t3)      
  
  return (aplicaSubst s3 t3, s3 ++ s2 ++ s1, n2 + 1)
\end{lstlisting}

El funcionament d'aquest bloc requereix quatre fases sincronitzades:
\begin{enumerate}
    \item S'infereix el tipus del descriptor esquerre (\texttt{e1}), generant un tipus presumptament funcional \texttt{t1} i les seves respectives substitucions \texttt{s1}.
    \item S'actualitza l'entorn de treball aplicant l'historial \texttt{s1} al context complet mitjançant \texttt{aplicaSubstCtx}, i a continuació s'infereix el tipus de l'argument dret (\texttt{e2}), obtenint \texttt{t2} i \texttt{s2}.
    \item S'encunya una variable fresca completament aliena \texttt{t3} per representar la naturalesa de retorn encara no descoberta de l'aplicació subjacent.
    \item Es propicia un escenari d'unificació: l'intèrpret força que el tipus esquerre (ja modelat prèviament i actualitzat amb el subconjunt \texttt{s2}) assumeixi formalment la identitat d'una funció \texttt{TFun t2 t3}. Posteriorment, el bloc retorna les capes de substitucions correctament combinades de dreta a esquerra (\texttt{s3 ++ s2 ++ s1}).
\end{enumerate}

\subsubsection{Estructures de Control i Definicions Recursives}

L'anàlisi del constructe \texttt{If e1 e2 e3} exigeix una coherència de dades transversal, garantint que la condició sigui booleana i que ambdues branques convergeixin unívocament cap a la mateixa essència de dada, generant diversos \texttt{sAcc} (acumuladors lineals de substitucions) per garantir que la informació teòrica s'escampi per tota l'expressió seqüencialment.

Paral·lelament, l'algorisme contempla el repte fonamental de les declaracions pròpies amb referències auto-contingudes. La funció auxiliar \texttt{infereixDefRec} està dissenyada explícitament per tractar funcions recursives introduint una suposició anticipada (una variable \texttt{t} fresca) dins del context lèxic sota el nom de la pròpia funció, permetent referenciar-la dins del seu cos abans d'inferir completament quin tipus definitiu representa. 

\section{Representació i Estètica Formats (Pretty-Printing)}

Més enllà del nucli de càlcul logicodeductiu intern, l'assoliment pràctic de tot aquest desplegament és lliurar al desenvolupador l'anàlisi semàntic de manera comprensible. Tot i que les variables formades numèricament (com \texttt{TVar "a1"}) proporcionen estabilitat algebraica per al processador HM, resulten redundants i confuses a l'hora de mostrar els resultats de tipus al bucle REPL de l'usuari \cite{hutton1998monadic}.

Per suplir aquesta mancança visual, s'ha implementat un mòdul de transformació estètica a la funció \texttt{showTipus}:

\begin{lstlisting}[style=haskellstyle, caption=Funció de representació textual amb respecte a l'associativitat per la dreta]
showTipus :: Tipus -> String
showTipus TInt = "Int"
showTipus TBool = "Bool"
showTipus (TFun (TFun t1 t2) t3) = "(" ++ showTipus t1 ++ " -> " ++ showTipus t2 ++ ")" ++ " -> " ++ showTipus t3 
showTipus (TFun t1 t2) = showTipus t1 ++ " -> " ++ showTipus t2
showTipus (TVar a) = a
\end{lstlisting}

Donat que en la teoria de tipus funcional el constructor de funcions és estructuralment \textbf{associatiu per la dreta}, una expressió visualment modelada com a $A \rightarrow B \rightarrow C$ ja implica intrínsecament l'ordre agrupacional correcte $A \rightarrow (B \rightarrow C)$ sense necessitar parèntesis. Ara bé, si ens trobem en el cas oposat on el paràmetre que es rep per la banda esquerra és ell mateix una altra funció d'ordre superior com $(A \rightarrow B) \rightarrow C$, els delimitadors de lectura prenen una funció cabdal de la gramàtica. 

La funció \texttt{showTipus} garanteix que la màquina discrimina de forma precisa quan un paràmetre esquerre posseeix forma funcional recursiva executant la branca \texttt{(TFun (TFun t1 t2) t3)}. Aquest factor aïlla la subestructura amb l'embolcall correcte per imprimir a la línia de comandes els tipus formals en el seu hàbit d'expressió canònica i fidel als cànons imposats a l'ecosistema Haskell clàssic.